\documentclass[12pt]{article}
\usepackage[utf8]{inputenc}
\usepackage[T1]{fontenc}
\usepackage{geometry}
\geometry{a4paper, margin=1in}
\usepackage{hyperref}
\usepackage{tabularx}
\usepackage{booktabs}
\usepackage{enumitem}
\usepackage{titlesec}

\titleformat{\section}{\large\bfseries}{\thesection}{1em}{}
\titleformat{\subsection}{\normalsize\bfseries}{\thesubsection}{1em}{}

\title{
    \huge \textbf{Assignment 7: Delivering a valuable product} \\
    \vspace{0.5cm}
    \large \textit{Poetry Recommender: Final Handover \& Transition}
}
\author{
    \textbf{Solo Developer: Zakhar} \\
    GitHub: \texttt{Potter32111}
}
\date{\today}

\begin{document}

\maketitle

\section*{Title Page Information}
\begin{itemize}
    \item \textbf{Project:} Poetry Recommender
    \item \textbf{Customer:} [TODO: Insert Customer Name]
    \item \textbf{Team Members:} Zakhar (Solo Founder)
    \item \textbf{Who did what:} End-to-end development, deployment, and transition documentation.
    \item \textbf{Who DID NOT participate:} N/A
\end{itemize}

\newpage

\section{Usefulness and Transition Report}

The final transition meeting for the \textbf{Poetry Recommender} project was held on March 22, 2026. The objective was to verify the delivery of MVP v3 and ensure the product is ready for long-term customer ownership.

\subsection{Product Completeness}
The product is considered feature-complete. All core functionalities defined in the MVP roadmap have been implemented and verified:
\begin{itemize}
    \item \textbf{Conversational Bot:} Robust Telegram interface using \texttt{aiogram}.
    \item \textbf{Recitation Check:} Voice-to-text validation using the \texttt{Vosk} engine.
    \item \textbf{Memory Retention:} SM-2 algorithm integration for spaced repetition.
    \item \textbf{Semantic Discovery:} \texttt{pgvector} based similarity search for poem recommendations.
\end{itemize}

\subsection{Customer Usage \& Feedback}
The customer has been using the product daily for personal poetry practice. Feedback has been overwhelmingly positive regarding the "silent" nature of the STT engine (offline processing) and the relevance of recommendations. The product is currently deployed on a \textbf{Hetzner VPS} and is fully accessible.

\subsection{Transition Strategy}
The transition involved a full handover of:
\begin{enumerate}
    \item \textbf{Source Code:} Clean, documented repository on GitHub.
    \item \textbf{Infrastructure:} Docker-compose setup for one-command deployment.
    \item \textbf{Maintenance:} Explicit instructions for database migrations and audio troubleshooting.
\end{enumerate}

\section{Customer Feedback on README}
During the handover, the README was audited with the client. While the installation steps were clear, the client requested two additional sections to empower self-sufficiency:
\begin{enumerate}
    \item \textbf{Troubleshooting Audio:} To handle device-specific voice recording issues.
    \item \textbf{Database Migrations:} To allow the client to update the schema as they add more custom content.
\end{enumerate}
These have been implemented as \texttt{### (for customer)} sections in the root \texttt{README.md}.

\section{Project Links}
The following links provide a comprehensive view of the final state of the project:
\begin{itemize}
    \item \textbf{Groomed Product Backlog:} [TODO: Insert Link]
    \item \textbf{Sprint Milestone (Final):} [TODO: Insert Link]
    \item \textbf{Sprint Tracking Table:} [TODO: Insert Link]
    \item \textbf{Extended MVP Roadmap:} [TODO: Insert Link]
    \item \textbf{Video Rehearsal (Demo Day):} [TODO: Insert Link]
\end{itemize}

\section{AI Usage Summary}
AI agents and LLMs were used throughout the 7 assignments for brainstorming, architecture drafting, test generation, and documentation refinement. A detailed log is available in \texttt{docs/reports/ai-usage.md}.

\end{document}
